\documentclass[11pt,a4paper]{article}
\usepackage[utf8]{inputenc}
\usepackage[T1]{fontenc}
\usepackage[spanish]{babel}
\usepackage[a4paper, margin=2.5cm]{geometry}
\usepackage{graphicx}
\usepackage{xcolor}
\usepackage{titlesec}
\usepackage{enumitem}
\usepackage{hyperref}
\usepackage{tikz}
\usepackage{booktabs}

\definecolor{velvetblack}{HTML}{0A0A0A}
\definecolor{velvetrose}{HTML}{B76E79}
\definecolor{velvetmauve}{HTML}{2B1F2A}
\definecolor{velvetchampagne}{HTML}{F4EADE}

\titleformat{\section}{\Large\bfseries\color{velvetrose}}{\thesection.}{0.5em}{}
\titleformat{\subsection}{\large\bfseries\color{velvetmauve}}{\thesubsection.}{0.5em}{}

\hypersetup{
  colorlinks=true,
  linkcolor=velvetrose,
  urlcolor=velvetrose,
  pdftitle={VELVET - Memoria descriptiva para registro de marca y app},
  pdfauthor={Sandro Siliuto}
}

\begin{document}

\begin{titlepage}
\centering
\vspace*{2cm}

\begin{tikzpicture}
  \draw[line width=1pt, velvetrose] (0,0) circle (1.2cm);
  \node at (0,0) {\Huge\color{velvetrose}\textbf{V}};
\end{tikzpicture}

\vspace{1.5cm}

{\Huge\bfseries\color{velvetblack} VELVET}\\[0.3cm]
{\Large\color{velvetrose} CONTACTOS}

\vspace{0.8cm}

{\large\textit{``EN LA VIDA TODO SON CONTACTOS''}}

\vspace{2cm}

{\LARGE\bfseries Memoria descriptiva para registro}\\[0.3cm]
{\LARGE\bfseries de marca y aplicación móvil}

\vfill

{\large Autor: \textbf{Sandro Siliuto}}\\[0.3cm]
{\large Fecha: \textbf{30 de junio de 2026}}

\end{titlepage}

\tableofcontents
\newpage

\section{Resumen ejecutivo}

VELVET es una aplicación móvil y plataforma digital de \textbf{networking social exclusivo} orientada a adultos que buscan establecer contactos de calidad en entornos de ocio nocturno, eventos privados y experiencias VIP. La propuesta de valor se centra en combinar un sistema de emparejamiento por tarjetas (swipe) con un ecosistema de recompensas geolocalizadas (VELVET\_GO), ofreciendo a los usuarios acceso a descuentos, consumiciones gratuitas, entradas VIP y merchandising exclusivo en locales asociados.

La identidad visual se inspira en la estética de un club privado nocturno: fondos oscuros, acentos en oro rosa, tipografías serif elegantes y efectos de vidrio esmerilado (glassmorphism). La palabra ``VIP'' está presente de forma recurrente en la interfaz para reforzar el posicionamiento de exclusividad.

\section{Concepto y propuesta de valor}

\subsection{Problema que resuelve}

En los entornos de ocio nocturno y eventos sociales, muchas personas encuentran dificultades para iniciar conversaciones o conocer a otras personas con intereses compatibles. Las aplicaciones de dating tradicionales carecen del contexto de evento en vivo, mientras que la experiencia presencial puede resultar intimidante o poco eficiente.

VELVET resuelve este problema mediante:

\begin{itemize}[leftmargin=*]
  \item \textbf{Perfiles verificados por teléfono}: cada usuario se registra con nombre y número de teléfono, reduciendo cuentas falsas.
  \item \textbf{Sistema de swipe en el entorno del evento}: los usuarios visualizan perfiles de personas presentes o cercanas y solo cuando ambas partes coinciden (match) se comparten los datos de contacto.
  \item \textbf{Geolocalización gamificada}: a través del módulo VELVET\_GO, los usuarios descubren checkpoints y recompensas en locales asociados, fomentando la interacción física.
\end{itemize}

\subsection{Público objetivo}

\begin{itemize}[leftmargin=*]
  \item Adultos de 25 a 45 años residentes en grandes ciudades españolas.
  \item Asistentes a eventos nocturnos, fiestas privadas, discotecas y festivales.
  \item Usuarios de aplicaciones de contactos que valoran la exclusividad y la seguridad.
\end{itemize}

\section{Funcionalidades principales}

\subsection{Registro y acceso VIP}

El acceso a VELVET comienza en una pantalla de bienvenida con el isotipo y la frase ``EN LA VIDA TODO SON CONTACTOS''. El usuario introduce:

\begin{itemize}[leftmargin=*]
  \item Nombre o apodo.
  \item Número de teléfono español (9 dígitos, empieza por 6 o 7).
  \item Foto opcional (solo se comparte si hay match).
\end{itemize}

El sistema almacena el perfil en una base de datos segura y genera una cookie de sesión para mantener la experiencia fluida. Si el usuario cierra sesión, puede volver a entrar con el mismo número de teléfono.

\subsection{Descubrimiento de perfiles (swipe)}

Una vez dentro, el usuario accede a una interfaz de tarjetas donde puede:

\begin{itemize}[leftmargin=*]
  \item Ver foto, nombre y datos de otros asistentes.
  \item Deslizar a la derecha (me gusta) o a la izquierda (descartar).
  \item Recibir notificación de match cuando ambos usuarios se gustan mutuamente.
  \item Acceder a la lista de matches y revisar perfiles conectados.
\end{itemize}

\subsection{Módulo VELVET\_GO: recompensas y checkpoints}

VELVET\_GO es la extensión gamificada de la aplicación. Funciona de forma similar a las geocompensas de Pokémon GO, pero adaptado al ocio nocturno:

\begin{itemize}[leftmargin=*]
  \item \textbf{Mapa de checkpoints}: muestra paradas exclusivas en locales asociados (discotecas, salas, rooftops).
  \item \textbf{Recompensas disponibles}: vales de descuento, copas gratis, entradas VIP y regalos o merchandising.
  \item \textbf{Desbloqueo por proximidad}: cuando el usuario está físicamente cerca de un checkpoint (dentro del radio configurado), puede desbloquear la recompensa.
  \item \textbf{Código de canje QR}: cada recompensa desbloqueada genera un código único (por ejemplo, VELVET-COPA-001) que el usuario presenta en el local para canjearla.
\end{itemize}

\subsection{Panel de administración}

Los organizadores y administradores de VELVET disponen de un panel web protegido para:

\begin{itemize}[leftmargin=*]
  \item Gestionar usuarios registrados.
  \item Crear, editar y eliminar recompensas.
  \item Crear, editar y eliminar checkpoints en el mapa.
  \item Controlar stock y fechas de validez de cada recompensa.
\end{itemize}

\section{Identidad de marca y sistema de diseño}

\subsection{Logotipo e isotipo}

El isotipo de VELVET representa una fusión simétrica de dos rostros humanos enfrentados, sugiriendo conexión, intimidad y emparejamiento selectivo. En el centro destaca una ``V'' que ancla el nombre de la aplicación y un destello de cuatro puntas que simboliza la chispa del encuentro y el estatus de élite.

\subsection{Paleta cromática}

\begin{center}
\begin{tabular}{llll}
\toprule
\textbf{Token} & \textbf{Nombre} & \textbf{HEX} & \textbf{Aplicación} \\
\midrule
Onyx Black & \#0A0A0A & RGB(10,10,10) & Fondo base \\
Velvet Mauve & \#2B1F2A & RGB(43,31,42) & Tarjetas y modales \\
Classic Rose Gold & \#B76E79 & RGB(183,110,121) & Botones y acentos \\
Deep Rose Gold & \#8F404C & RGB(143,64,76) & Bordes y estados de presión \\
Pale Rose Gold & \#F2D7D3 & RGB(242,215,211) & Etiquetas y textos secundarios \\
Brut Champagne & \#F4EADE & RGB(244,234,222) & Texto principal \\
\bottomrule
\end{tabular}
\end{center}

\subsection{Tipografía}

\begin{itemize}[leftmargin=*]
  \item \textbf{Títulos editoriales}: Cinzel (serif elegante, estilo moda y alta joyería).
  \item \textbf{Texto de interfaz y cuerpo}: Inter (sans-serif geométrica, alta legibilidad).
\end{itemize}

\subsection{Estilo visual}

La interfaz emplea \textbf{glassmorphism}: fondos translúcidos con desenfoque, bordes sutiles en tonos oro rosa y sombras difusas que generan sensación de profundidad y lujo. Los botones principales utilizan un gradiente metálico de oro rosa que emula reflejos de metal pulido.

\section{Arquitectura técnica}

\subsection{Stack tecnológico}

\begin{itemize}[leftmargin=*]
  \item \textbf{Frontend}: Next.js 15 + React + TypeScript + Tailwind CSS.
  \item \textbf{Backend}: API routes serverless de Next.js.
  \item \textbf{Base de datos}: Supabase (PostgreSQL).
  \item \textbf{Almacenamiento}: Supabase Storage para fotos de perfil y recompensas.
  \item \textbf{Mapas}: Leaflet con capas oscuras de CARTO.
  \item \textbf{Despliegue}: Vercel.
\end{itemize}

\subsection{Tablas principales en Supabase}

\begin{itemize}[leftmargin=*]
  \item \texttt{velvet\_users}: perfiles de usuario (id, nombre, teléfono, foto).
  \item \texttt{swipes}: acciones de like/dislike entre usuarios.
  \item \texttt{matches}: coincidencias mutuas.
  \item \texttt{rewards}: catálogo de recompensas VIP.
  \item \texttt{checkpoints}: ubicaciones geográficas que desbloquean recompensas.
  \item \texttt{user\_rewards}: recompensas desbloqueadas por cada usuario.
  \item \texttt{checkpoint\_visits}: registro de visitas a checkpoints.
\end{itemize}

\subsection{Seguridad y privacidad}

\begin{itemize}[leftmargin=*]
  \item Los números de teléfono son únicos para evitar duplicados.
  \item Las fotos y datos de contacto solo se comparten cuando existe match mutuo.
  \item Las operaciones de escritura en la base de datos se realizan desde el servidor mediante claves de servicio seguras.
  \item Row Level Security (RLS) en Supabase para controlar el acceso a datos.
\end{itemize}

\section{URLs de despliegue y repositorios}

\begin{itemize}[leftmargin=*]
  \item \textbf{VELVET contactos}: \url{https://vamos-a-definir-la-verbena-clone.vercel.app}
  \item \textbf{VELVET\_GO}: \url{https://vamos-a-definir-la-go.vercel.app}
  \item \textbf{Repositorio VELVET\_GO}: \url{https://github.com/sandrosiliuto/velvet-go}
\end{itemize}

\section{Clases de productos y servicios sugeridas para registro}

Para el registro de marca y/o aplicación, se sugieren las siguientes clases del Nice:

\begin{itemize}[leftmargin=*]
  \item \textbf{Clase 9}: software descargable para contactos sociales, emparejamiento y gestión de recompensas; aplicaciones móviles.
  \item \textbf{Clase 35}: servicios de publicidad, promoción de eventos y fidelización de clientes mediante recompensas.
  \item \textbf{Clase 38}: servicios de telecomunicaciones y mensajería entre usuarios.
  \item \textbf{Clase 41}: organización de eventos sociales, fiestas y ocio nocturno.
  \item \textbf{Clase 42}: desarrollo y alojamiento de software y plataformas digitales.
\end{itemize}

\section{Declaración de autoría}

La presente memoria describe una aplicación y marca original concebida, desarrollada y desplegada por \textbf{Sandro Siliuto}. El código fuente, el diseño de interfaz, la identidad visual y la lógica de negocio han sido creados de forma original para proteger la idea, el software y la marca comercial asociada a VELVET.

\end{document}
